% ================================================================
%  sesion14.tex  —  Sesión 14 de 20
%  Leer e interpretar gráficas
%  Bloque 3: Datos y gráficas · Matemáticas 1.º ESO
% ================================================================
\documentclass[aspectratio=169, 12pt, spanish]{beamer}
\usepackage{almeria-beamer}

\renewcommand{\SessionNumber}{14}
\renewcommand{\SessionTitle}{Leer e interpretar gráficas}
\renewcommand{\SessionName}{Sesión 14}
\renewcommand{\SessionObjective}{%
  Leer coordenadas de una gráfica, identificar tendencias,
  máximos, mínimos y valores clave, y extraer conclusiones
  sobre la realidad de Almería.}
\renewcommand{\BlockName}{Bloque 3: Datos y gráficas}

\begin{document}

\begin{frame}[plain]\titlepage\end{frame}

\begin{frame}{Índice}
  \begin{columns}[t]
    \begin{column}{0.5\textwidth}
      \begin{enumerate}
        \item Cómo leer un valor de una gráfica
        \item Elementos clave: máximo, mínimo, tendencia
        \item Interpretación de la pendiente
        \item Gráfica de temperatura en Almería
        \item Gráfica de visitantes 2017--2022
      \end{enumerate}
    \end{column}
    \begin{column}{0.5\textwidth}
      \begin{enumerate}\setcounter{enumi}{5}
        \item Preguntas tipo sobre gráficas
        \item Ejercicios multinivel (Bloom)
        \item Evidencia del día
        \item Error frecuente
        \item Resumen
      \end{enumerate}
    \end{column}
  \end{columns}
\end{frame}

% ---------------------------------------------------------------
\seccionframe{1. Cómo leer un valor de una gráfica}
% ---------------------------------------------------------------

\begin{frame}{Protocolo de lectura de una gráfica}
  \begin{columns}[c]
    \begin{column}{0.5\textwidth}
      \begin{pasobox}[1]
        \footnotesize
        Localiza el valor de \textbf{$x$} que te interesa
        en el eje horizontal.
      \end{pasobox}
      \vspace{4pt}
      \begin{pasobox}[2]
        \footnotesize
        Sube en vertical hasta tocar la
        \textbf{curva o barra}.
      \end{pasobox}
      \vspace{4pt}
      \begin{pasobox}[3]
        \footnotesize
        Desplázate en horizontal hacia el
        \textbf{eje vertical} y lee el valor $y$.
      \end{pasobox}
      \vspace{4pt}
      \begin{pasobox}[4]
        \footnotesize
        Escribe el resultado con sus
        \textbf{unidades} (no olvides el eje $y$).
      \end{pasobox}
    \end{column}
    \begin{column}{0.46\textwidth}
      \begin{tikzpicture}[scale=0.82]
        \begin{axis}[
          almeria line,
          width=6.8cm, height=5.2cm,
          xmin=0, xmax=7, ymin=0, ymax=30,
          xtick={1,2,3,4,5,6},
          xticklabels={Ene,Feb,Mar,Abr,May,Jun},
          ytick={0,10,20,30},
          ylabel={Temperatura (°C)},
          xlabel={Mes},
        ]
          \addplot[AlmTerra, line width=2pt, mark=*]
            coordinates {(1,12)(2,13)(3,16)(4,19)(5,23)(6,27)};
          % Líneas de lectura para Abril (x=4, y=19)
          \draw[AlmVerde, dashed, line width=1pt]
            (axis cs:4,0) -- (axis cs:4,19) -- (axis cs:0,19);
          \node[fill=AlmVerde!20, draw=AlmVerde, rounded corners,
                font=\tiny\bfseries, inner sep=2pt]
            at (axis cs:2.5,19) {$y=19\,°C$};
          \node[fill=AlmVerde!20, draw=AlmVerde, rounded corners,
                font=\tiny\bfseries, inner sep=2pt]
            at (axis cs:4,5) {$x=\text{Abr}$};
        \end{axis}
      \end{tikzpicture}
    \end{column}
  \end{columns}
\end{frame}

% ---------------------------------------------------------------
\seccionframe{2. Elementos clave de una gráfica}
% ---------------------------------------------------------------

\begin{frame}{Máximo, mínimo, tendencia y punto de inflexión}
  \begin{columns}[T]
    \begin{column}{0.42\textwidth}
      \begin{definicionbox}{Valor máximo}
        Punto más alto de la gráfica.
        Mayor valor de $y$.
      \end{definicionbox}
      \vspace{5pt}
      \begin{definicionbox}{Valor mínimo}
        Punto más bajo de la gráfica.
        Menor valor de $y$.
      \end{definicionbox}
      \vspace{5pt}
      \begin{definicionbox}{Tendencia}
        Dirección general de la curva:\\
        $\nearrow$ creciente \quad
        $\searrow$ decreciente \quad
        $\rightarrow$ constante
      \end{definicionbox}
      \vspace{5pt}
      \begin{definicionbox}{Punto de inflexión}
        Donde la tendencia cambia de sentido.
      \end{definicionbox}
    \end{column}
    \begin{column}{0.54\textwidth}
      \begin{tikzpicture}
        \begin{axis}[
          almeria line,
          width=8.5cm, height=5.5cm,
          xmin=1, xmax=12, ymin=10, ymax=33,
          xtick={1,2,3,4,5,6,7,8,9,10,11,12},
          xticklabels={E,F,M,A,M,J,J,A,S,O,N,D},
          ylabel={°C}, xlabel={Mes},
          title={Temperatura media mensual en Almería},
        ]
          \addplot[AlmTerra, line width=2pt, mark=*,
                   mark options={fill=AlmTerra}]
            coordinates {
              (1,12)(2,13)(3,16)(4,19)(5,23)(6,27)
              (7,30)(8,30)(9,26)(10,21)(11,16)(12,13)
            };
          % Máximo
          \node[above, font=\tiny\bfseries, color=AlmTerra]
            at (axis cs:7.5,30) {\faArrowUp\ Máx: 30°C};
          % Mínimo
          \node[below right, font=\tiny\bfseries, color=AlmAzulM]
            at (axis cs:1,12) {\faArrowDown\ Mín: 12°C};
          % Tendencia creciente
          \draw[->, AlmVerde, line width=1.2pt]
            (axis cs:2,14) -- (axis cs:6,26)
            node[midway, above left, font=\tiny, color=AlmVerde]
            {creciente};
          % Tendencia decreciente
          \draw[->, BEvaluar, line width=1.2pt]
            (axis cs:8,29) -- (axis cs:11,17)
            node[midway, above right, font=\tiny, color=BEvaluar]
            {decreciente};
        \end{axis}
      \end{tikzpicture}
    \end{column}
  \end{columns}
\end{frame}

% ---------------------------------------------------------------
\seccionframe{3. Interpretación de la pendiente}
% ---------------------------------------------------------------

\begin{frame}{¿Qué nos dice la pendiente de una curva?}
  \begin{columns}[c]
    \begin{column}{0.5\textwidth}
      \begin{teoriabox}
        La \textbf{pendiente} (inclinación) indica la
        \textit{velocidad de cambio}:
        \begin{itemize}
          \item Pendiente \textbf{pronunciada}: el valor
                cambia \textbf{rápido}.
          \item Pendiente \textbf{suave}: el valor cambia
                \textbf{lentamente}.
          \item Línea \textbf{horizontal}: el valor
                \textbf{no cambia}.
        \end{itemize}
      \end{teoriabox}
    \end{column}
    \begin{column}{0.46\textwidth}
      \begin{tikzpicture}[scale=0.85]
        \begin{axis}[
          almeria line,
          width=6.8cm, height=4.8cm,
          xmin=0, xmax=6, ymin=0, ymax=12,
          axis lines=left,
          xtick={1,2,3,4,5},
          ytick={0,4,8,12},
          ylabel={$y$}, xlabel={$x$},
        ]
          % Pendiente pronunciada
          \addplot[AlmTerra, line width=2pt]
            coordinates {(0,0)(1,4)(2,8)(3,12)};
          \node[right, font=\tiny\bfseries, color=AlmTerra]
            at (axis cs:3,12) {pronunciada};
          % Pendiente suave
          \addplot[AlmVerde, line width=2pt]
            coordinates {(0,0)(1,1.5)(2,3)(3,4.5)(5,7.5)};
          \node[right, font=\tiny\bfseries, color=AlmVerde]
            at (axis cs:5,7.5) {suave};
          % Horizontal
          \addplot[AlmAzulM, line width=2pt, dashed]
            coordinates {(0,2)(5,2)};
          \node[right, font=\tiny\bfseries, color=AlmAzulM]
            at (axis cs:4,2.5) {constante};
        \end{axis}
      \end{tikzpicture}
    \end{column}
  \end{columns}
\end{frame}

% ---------------------------------------------------------------
\seccionframe{4. Temperatura en Almería}
% ---------------------------------------------------------------

\begin{frame}{Análisis de la gráfica de temperatura}
  \begin{columns}[c]
    \begin{column}{0.55\textwidth}
      \begin{tikzpicture}
        \begin{axis}[
          almeria line,
          width=9cm, height=5.5cm,
          xmin=1, xmax=12, ymin=8, ymax=34,
          xtick={1,...,12},
          xticklabels={E,F,M,A,M,J,J,A,S,O,N,D},
          ylabel={Temperatura (°C)},
          xlabel={Mes},
          title={Temperatura media · Almería},
          ytick={10,15,20,25,30},
        ]
          \addplot[AlmTerra, line width=2.5pt, mark=*,
                   mark options={fill=AlmTerra, scale=1.2}]
            coordinates {
              (1,12)(2,13)(3,16)(4,19)(5,23)(6,27)
              (7,30)(8,30)(9,26)(10,21)(11,16)(12,13)
            };
          % Línea de referencia 20°C
          \addplot[AlmAmbar!70, dashed, line width=1pt]
            coordinates {(1,20)(12,20)}
            node[right, font=\tiny, color=AlmAmbar!80!black]
            {20°C};
        \end{axis}
      \end{tikzpicture}
    \end{column}
    \begin{column}{0.42\textwidth}
      Respondemos las preguntas clave:
      \begin{itemize}
        \item \textbf{Máximo:} julio--agosto, $30\,°C$
        \item \textbf{Mínimo:} enero, $12\,°C$
        \item \textbf{Rango:} $30-12 = 18\,°C$
        \item \textbf{Meses $>20\,°C$:} junio, julio, agosto, septiembre
        \item \textbf{Tendencia} ene→ago: \textit{creciente}
        \item \textbf{Tendencia} ago→dic: \textit{decreciente}
        \item \textbf{Simetría:} casi simétrica respecto a agosto
      \end{itemize}
      \vspace{6pt}
      \begin{notabox}
        \footnotesize
        Almería tiene uno de los climas
        más cálidos y secos de Europa.
        ¡Más de 3\,000 horas de sol al año!
      \end{notabox}
    \end{column}
  \end{columns}
\end{frame}

% ---------------------------------------------------------------
\seccionframe{5. Visitantes 2017--2022}
% ---------------------------------------------------------------

\begin{frame}{Turismo en Almería: gráfica de barras 2017--2022}
  \begin{columns}[c]
    \begin{column}{0.58\textwidth}
      \begin{tikzpicture}
        \begin{axis}[
          almeria bar,
          width=9.5cm, height=5.5cm,
          symbolic x coords={2017,2018,2019,2020,2021,2022},
          xtick=data,
          ylabel={Visitantes (millones)},
          xlabel={Año},
          title={Visitantes a Almería (millones)},
          ymin=0, ymax=3.5,
          nodes near coords style={font=\tiny\bfseries},
          bar width=0.55cm,
        ]
          \addplot[fill=AlmAzulM, fill opacity=0.85]
            coordinates {
              (2017,2.5)(2018,2.7)(2019,2.9)
              (2020,1.2)(2021,1.8)(2022,2.8)
            };
          % Línea de tendencia general
          \addplot[AlmTerra, line width=1.5pt, dashed, no marks]
            coordinates {(2017,2.5)(2018,2.7)(2019,2.9)
                          (2020,1.2)(2021,1.8)(2022,2.8)};
        \end{axis}
      \end{tikzpicture}
    \end{column}
    \begin{column}{0.38\textwidth}
      Datos relevantes:
      \begin{itemize}
        \item Máximo: 2019 → 2,9 M
        \item Mínimo: 2020 → 1,2 M
        \item Caída 2019→2020:\\
          $\dfrac{2{,}9-1{,}2}{2{,}9}\times 100
          \approx \mathbf{59\,\%}$
        \item Recuperación 2021→2022:\\
          $\dfrac{2{,}8-1{,}8}{1{,}8}\times 100
          \approx \mathbf{56\,\%}$
      \end{itemize}
      \vspace{6pt}
      \begin{analizarbox}
        \footnotesize
        ¿Qué evento de 2020 explica
        la caída drástica de visitantes?
        ¿Qué nos indica la recuperación
        de 2022?
      \end{analizarbox}
    \end{column}
  \end{columns}
\end{frame}

% ---------------------------------------------------------------
\seccionframe{6. Preguntas tipo sobre gráficas}
% ---------------------------------------------------------------

\begin{frame}{Preguntas que siempre debes poder responder}
  \begin{columns}[T]
    \begin{column}{0.48\textwidth}
      \textbf{Lectura de valores:}
      \begin{itemize}
        \item «¿Cuál es el valor de $y$ cuando $x = \ldots$?»
        \item «¿En qué mes/año fue el valor más alto?»
        \item «¿Entre qué valores oscila $y$?»
      \end{itemize}
      \vspace{8pt}
      \textbf{Tendencias:}
      \begin{itemize}
        \item «¿La gráfica es creciente, decreciente
              o constante entre $x_1$ y $x_2$?»
        \item «¿Dónde está el punto de inflexión?»
        \item «¿Cuál es la tendencia general?»
      \end{itemize}
    \end{column}
    \begin{column}{0.48\textwidth}
      \textbf{Comparación:}
      \begin{itemize}
        \item «¿Entre qué dos puntos fue mayor
              el crecimiento?»
        \item «¿Cuál es la variación total de $y$?»
      \end{itemize}
      \vspace{8pt}
      \textbf{Interpretación:}
      \begin{itemize}
        \item «¿Qué puede explicar la anomalía en $x = \ldots$?»
        \item «¿Qué predices para el siguiente valor?»
        \item «¿Qué conclusión puedes extraer del gráfico?»
      \end{itemize}
      \vspace{8pt}
      \begin{notabox}[Consejo]
        \footnotesize
        Siempre lee el \textbf{título}, los
        \textbf{ejes} y la \textbf{unidad} antes
        de responder cualquier pregunta.
      \end{notabox}
    \end{column}
  \end{columns}
\end{frame}

% ---------------------------------------------------------------
\seccionframe{7. Ejercicios multinivel}
% ---------------------------------------------------------------

\begin{frame}{Ejercicios multinivel}
  \begin{columns}[T]
    \begin{column}{0.31\textwidth}
      \begin{comprenderbox}
        \footnotesize\dificultad{1}\\[4pt]
        Usando la gráfica de temperatura de Almería:\\[4pt]
        (a) ¿Temperatura en abril?\\
        (b) ¿Mes más caluroso?\\
        (c) ¿Cuántos meses supera 20°C?\\
        (d) Describe la tendencia de enero a agosto
        en una frase.
      \end{comprenderbox}
    \end{column}
    \begin{column}{0.31\textwidth}
      \begin{aplicarbox}
        \footnotesize\dificultad{2}\\[4pt]
        Con la gráfica de visitantes:\\[4pt]
        (a) Mejor año y número de visitantes.\\
        (b) Peor año y número de visitantes.\\
        (c) Calcula el \% de cambio de 2021 a 2022.\\
        (d) Si la tendencia continúa, ¿cuántos
        visitantes habría en 2023?
      \end{aplicarbox}
    \end{column}
    \begin{column}{0.31\textwidth}
      \begin{analizarbox}
        \footnotesize\dificultad{3}\\[4pt]
        La gráfica muestra el número de terrazas
        abiertas en Almería por mes: máximo en
        agosto, mínimo en enero.\\[4pt]
        \textbf{Analiza:} ¿Qué relación tiene con
        la temperatura? ¿Con el turismo?
        ¿Qué variable \textit{causa} la apertura
        de terrazas? Argumenta en 6 líneas.
      \end{analizarbox}
    \end{column}
  \end{columns}
\end{frame}

% ---------------------------------------------------------------
\begin{frame}{Evidencia del día}
  \begin{evidenciabox}
    Escribe un \textbf{análisis de dos gráficas} (temperatura y
    visitantes) que incluya:
    \begin{enumerate}
      \item Los valores leídos en al menos 3 puntos de cada gráfica
      \item Identificación del máximo, mínimo y tendencia general
      \item Una anomalía o punto de inflexión con su posible explicación
      \item Una conclusión de 3--4 líneas sobre lo que la gráfica revela
            sobre Almería
    \end{enumerate}
    \vspace{4pt}
    \textit{No basta con leer los números: debes interpretarlos.}
  \end{evidenciabox}
\end{frame}

% ---------------------------------------------------------------
\begin{frame}{¡Cuidado con este error!}
  \begin{errorbox}
    \begin{columns}[c]
      \begin{column}{0.48\textwidth}
        \incorrecto\ Leer los ejes al revés:\\[4pt]
        «En el mes de julio, la temperatura
        es 7» ← Leyó el eje $x$ donde debía
        leer el eje $y$.
      \end{column}
      \begin{column}{0.48\textwidth}
        \correcto\ Protocolo correcto:\\[4pt]
        1. Localiza julio en el eje \textbf{horizontal}.\\
        2. Sube hasta la curva.\\
        3. Lee el valor en el eje \textbf{vertical}: $30\,°C$.
      \end{column}
    \end{columns}
  \end{errorbox}
  \vspace{6pt}
  \begin{errorbox}
    \incorrecto\ «En 2020 bajaron los visitantes porque
    no había buenos hoteles.»\\[4pt]
    \correcto\ Una gráfica muestra \textbf{qué} pasó,
    pero no siempre explica \textbf{por qué}.
    Hay que buscar la explicación con otras fuentes o razonar
    con más datos (en 2020 hubo pandemia mundial de COVID-19).
  \end{errorbox}
\end{frame}

% ---------------------------------------------------------------
\begin{frame}{Resumen de la sesión 14}
  \begin{resumenbox}
    \begin{itemize}
      \item Para leer un valor: localiza $x$ → sube a la curva
            → lee $y$ en el eje vertical.
      \item Identificamos \textbf{máximo}, \textbf{mínimo},
            \textbf{tendencia} (creciente/decreciente/constante)
            y \textbf{punto de inflexión}.
      \item La pendiente indica la \textbf{velocidad de cambio}.
      \item La gráfica de temperatura de Almería es \textit{creciente}
            de enero a julio--agosto, y \textit{decreciente} de agosto
            a diciembre.
      \item Anomalías (como 2020) requieren \textbf{interpretación
            contextual}, no solo matemática.
    \end{itemize}
  \end{resumenbox}
  \vspace{6pt}
  \begin{center}
    {\small\color{AlmAzulM}\faArrowCircleRight\enskip
    \textbf{Sesión 15:} Comparar gráficas y detectar correlaciones}
  \end{center}
\end{frame}

\end{document}
